\usepackage{xeCJK}
\usepackage{geometry}
\geometry{left=2cm,right=2cm,top=2.5cm,bottom=2.5cm}
\setlength{\headheight}{12.7pt}

\usepackage{titlesec}
\titleformat{\section}
  {\normalfont\large\bfseries}
  {\thesection}
  {1em}
  {}
\titlespacing*{\section}
  {0pt}
  {3.5ex plus 1ex minus .2ex}
  {2.3ex plus .2ex}

\usepackage{amsmath}
\usepackage{amssymb}
\usepackage{amsthm}
\usepackage{enumitem}
\setlist[enumerate]{itemsep=0pt,topsep=0pt,parsep=0pt,leftmargin=0pt,itemindent=2.5em}
\setlist[itemize]{itemsep=0pt,topsep=0pt,parsep=0pt,leftmargin=0pt,itemindent=2.5em}
\usepackage{fancyhdr}
\pagestyle{fancy}
\fancyhf{}
\usepackage{graphicx}
\usepackage{hyperref}
\usepackage{indentfirst}
\usepackage{lmodern}


\begin{document}
\setlength{\abovedisplayskip}{1pt}
\setlength{\belowdisplayskip}{1pt}

\section{序数}

序数用于编号和排序.

\vspace{10pt}

\hypertarget{sec:定义1.1}{}
\noindent\textbf{定义1.1 (序数)}

给定集合$\alpha$. 如果$\alpha$
是\href{../../03 自然数/01 自然数/02 传递集#nameddest=sec:定义1.1}{传递集}并且
$(\alpha,\in)$是\href{../01 序关系/04 良序关系#nameddest=sec:定义1.1}{严格良序},
那么称$\alpha$是一个\textbf{序数}.

\vspace{10pt}

显然\href{../../03 自然数/01 自然数/01 自然数#nameddest=sec:定义1.3}{自然数}和
自然数集$\omega$都是序数.
根据正则公理, 序数有更简单的判定方法.

\vspace{10pt}

\hypertarget{sec:定理1.1}{}
\noindent\textbf{定理1.1}

任给集合$\alpha$, $\alpha$是序数当且仅当$\alpha$是传递集并且
$\forall x,y\in\alpha:x\in y\,\lor\,x=y\,\lor\,y\in x$.

\noindent{\kaishu 证明.}

\noindent{\kaishu 必要性.} 显然.

\noindent{\kaishu 充分性.}

假设$\alpha$是传递集, 并且$(\alpha,\in)$有三歧性. 此时$(\alpha,\in)$一定是严格良序:
\begin{enumerate}
    \item 反自反性: 由正则公理即得.
    \item 传递性: 在$\alpha$中如果$x\in y\in z$, 那么假设$x=z$或$z\in x$都有矛盾.
    由三歧性可知$x\in z$.
    \item 三歧性: 已知.
    \item 良基性: 由正则公理即得. \hfill $\qedsymbol$
\end{enumerate}

\vspace{10pt}

\hypertarget{sec:定理1.2}{}
\noindent\textbf{定理1.2}

\begin{enumerate}
    \item 任给序数$\alpha$和$\beta\in\alpha$, $\beta$也是序数,
    \item 任给序数$\alpha$和$\beta\in\alpha$, $\alpha_\beta=\beta$,
    \item 任给序数$\alpha,\beta$,
    如果$(\alpha,\in)\cong(\beta,\in)$, 那么$\alpha=\beta$,
    \item 任给序数$\alpha,\beta$,
    有$\alpha\in\beta\,\lor\,\alpha=\beta\,\lor\,\beta\in\alpha$,
    \item 任给序数$\alpha,\beta$,
    $\alpha\subset\beta$当且仅当$\alpha\in\beta\,\lor\,\alpha=\beta$.
\end{enumerate}

\noindent{\kaishu 证明.}

\begin{enumerate}
    \item 第一, 如果$x\in y\in\beta$, 那么$x,y\in\alpha$, 由传递性可知$x\in\beta$.
    所以$\beta$是传递集.
    
    \hspace{2.17em}
    第二, $\beta\subset\alpha$, 所以$(\beta,\in)$是严格良序.

    \item 对任意的$x$,
    $x\in\alpha_\beta\iff x\in\alpha\land x\in\beta\iff x\in\beta$.

    \item 取$f:(\alpha,\in)\cong(\beta,\in)$, 对$x\in\alpha$归纳证明$f(x)=x$.
    任取$x\in\alpha$, 假设$\forall y\in x:f(y)=y$.

    \hspace{2.17em}
    一方面, 对任意的$t\in f(x)$, 有$f^{-1}(t)\in x$,
    依假设有$t=f(f^{-1}(t))=f^{-1}(t)\in x$;

    \hspace{2.17em}
    另一方面, 对任意的$t\in x$, 依假设$t=f(t)\in f(x)$.

    \hspace{2.17em}
    综上$f(x)=x$, 归纳得$\forall x\in\alpha:f(x)=x$. 所以$\beta=\alpha$.

    \item 根据\href{../01 序关系/04 良序关系#nameddest=sec:定理2.4}{良序集基本定理},
    有以下的情况:

    \hspace{2.17em}
    如果$(\alpha,\in)\cong(\beta,\in)$, 那么$\alpha=\beta$;

    \hspace{2.17em}
    如果存在$x\in\alpha$使得$(\alpha_x,\in)\cong(\beta,\in)$, 那么
    $\beta=\alpha_x=x\in\alpha$;

    \hspace{2.17em}
    如果存在$y\in\beta$使得$(\alpha,\in)\cong(\beta_y,\in)$, 那么
    $\alpha=\beta_y=y\in\beta$.

    \item 如果$\alpha\subset\beta$, 假设$\beta\in\alpha$则$\beta\in\beta$矛盾,
    所以$\alpha\in\beta\,\lor\,\alpha=\beta$;

    \hspace{2.17em}
    如果$\alpha\in\beta\,\lor\,\alpha=\beta$,
    显然$\alpha\subset\beta$. \hfill $\qedsymbol$
\end{enumerate}





\newpage

接下来研究序数组成的集合.

\vspace{10pt}

\hypertarget{sec:定理1.3}{}
\noindent\textbf{定理1.3}

任给非空集合$A$, 如果$A$的元素都是序数, 那么:
\vspace{3pt}
\begin{enumerate}
    \item $(A,\in)$是严格良序,
    \vspace{3pt}
    \item $\displaystyle\bigcup A$是序数, 并且$\displaystyle\bigcup A=\sup A$,
    \vspace{3pt}
    \item $\displaystyle\bigcap A$是序数, 并且$\displaystyle\bigcap A=\min A$.
\end{enumerate}

\noindent{\kaishu 证明.}

\begin{enumerate}
    \item 传递性: 在$A$中如果$\alpha\in\beta\in\gamma$,
    那么由$\gamma$是传递集得$\alpha\in\beta$.

    \hspace{2.17em}
    三歧性: 在$A$中任取$\alpha,\beta$,
    有$\alpha\in\beta\,\lor\,\alpha=\beta\,\lor\,\beta\in\alpha$.

    \item 第一, $A$的元素都是传递集, 所以$\displaystyle\bigcup A$是传递集.
    
    \vspace{3pt}\hspace{2.17em}
    第二, 对任意的$\displaystyle x\in\bigcup A$, 存在$\alpha\in A$使得$x\in\alpha$,
    于是$x$是序数.
    
    \vspace{3pt}\hspace{2.17em}
    从而$\displaystyle\bigcup A$的元素都是序数,
    即$\displaystyle\left(\bigcup A,\in\right)$是严格良序.
    所以$\displaystyle\bigcup A$是序数.

    \vspace{3pt}\hspace{2.17em}
    进一步, 对任意的$\alpha\in A$有$\displaystyle\alpha\subset\bigcup A$,
    所以$\displaystyle\bigcup A$是$A$的上界;

    \vspace{3pt}\hspace{2.17em}
    如果$\beta$是更小的上界, 即$\displaystyle\beta\in\bigcup A$,
    那么存在$\alpha\in A$使得$\beta\in\alpha$, 矛盾.

    \vspace{3pt}
    \item 第一, $A$的元素都是传递集, 所以$\displaystyle\bigcap A$是传递集.
    
    \vspace{3pt}\hspace{2.17em}
    第二, $\displaystyle\bigcap A\subset\bigcup A$,
    从而$\displaystyle\bigcap A$的元素也都是序数,
    即$\displaystyle\left(\bigcap A,\in\right)$是严格良序.
    所以$\displaystyle\bigcap A$是序数.

    \vspace{3pt}\hspace{2.17em}
    进一步, 记$m=\min A$. 那么:

    \vspace{3pt}\hspace{2.17em}
    一方面, 对任意的$\displaystyle x\in\bigcap A$, 由$m\in A$即得$x\in m$;

    \vspace{3pt}\hspace{2.17em}
    另一方面, 对任意的$x\in m$, 此时对任意的$\alpha\in A$有$m\subset\alpha$,
    从而$x\in\alpha$, 即$\displaystyle x\in\bigcap A$.
    
    \vspace{3pt}\hspace{2.17em}
    综上$\displaystyle\bigcap A=m=\min A$. \hfill $\qedsymbol$
\end{enumerate}

\vspace{10pt}

以后可以把序数间的$\in$写成$<$, 序数间的$\subset$写成$\leqslant$.

\vspace{10pt}

\hypertarget{sec:定义1.2}{}
\noindent\textbf{定义1.2 (序数类)}

记全体序数组成的类为\textbf{序数类}, 记作$\mathbf{On}$.
序数类是\href{../../01 ZFC 公理/02 类#nameddest=sec:真类}{真类}.

\noindent{\kaishu 验证.}

假设存在集合$\mathrm{On}$使得
$\forall\alpha:\alpha\in\mathrm{On}\leftrightarrow\alpha\text{\,是序数\,}$. 那么:
\begin{enumerate}
    \item 序数的元素都是序数, 所以$\mathrm{On}$是传递集,
    \item $\mathrm{On}$的元素都是序数, 所以$(\mathrm{On},\in)$是严格良序,
\end{enumerate}
所以$\mathrm{On}$也是序数, 即$\mathrm{On}\in\mathrm{On}$, 矛盾. \hfill $\qedsymbol$

\end{document}